% Experimental section for the main AAMAS manuscript.
\providecommand{\AAMASExperimentFigureDir}{figures}

\section{Experimental Evaluation}
\label{sec:cage-experiments}
\label{sec:switching-experiments}

\paragraph{Observed demand shares.}
We replay 730 chronological days of NYC Forestry Service Requests in the
Hazard category, 2021--2022~\cite{nycforestry2026}. Daily counts are grouped
as Brooklyn, Queens, and the other three boroughs and normalized to
$\ell_t\in\Delta_3$; the single zero-count day uses uniform shares.
These administrative counts define a modeled allocation objective, rather
than measured staffing needs or service outcomes. For
$P=L=\Delta_3$, set
\[
 u(p,\ell)=(p-\ell)/\sqrt2,\qquad p^\star(\ell)=\ell.
\]
Thus $S(Q)=\{0\}$ for every nonempty $Q$, including all interior mixtures,
and $q_T\le2$.

\paragraph{Certified budget and controls.}
A fresh lag base starts uniformly at $p_0$ and then plays the previous
observed share. On any adaptive $h$-round continuation,
$\sum_{j=1}^h u(p_j,\ell_j)=(p_0-\ell_h)/\sqrt2$, whose norm is at most one.
Consequently $B_0(0)=0$, $B_0(h)=1$ for $h>0$, and $G_T=1$ are valid
before observing the trace. This instantiates Theorem~\ref{thm:master}
with a game-specific base; it is separate from the block routine of
Proposition~\ref{prop:safe}.
We compare the master, fast-only, lag-safe-only, a trailing-mean window
($W=16$), and block-safe-only. Every method starts uniformly, chooses
before observing that day's share, and runs continuously for the announced
$T=730$. The window and budgets are fixed, without tuning on this trace.

\paragraph{Metrics and results.}
We report $\delta_t=\|t^{-1}\sum_{s\le t}u_s\|_2$ and the additional daily
mismatch $M_T=T^{-1}\sum_{t\le T}\|u_t\|_2$. The former permits cancellation
of signed allocation imbalances; the latter measures daily tracking.
The master crosses its budget on day 44 (13 February 2021), with
$E_{44}=1.090736$, and uses the fresh lag base for the remaining 686 days.
Table~\ref{tab:switching-nyc} and Figure~\ref{fig:switching-nyc} show different
rankings: switching reduces daily mismatch relative to fast-only, while
fast-only has a smaller terminal $\delta_T$. Lag and window controls have
smaller daily mismatch than the master. These are exploratory fixed-trace
comparisons; no confidence intervals or statistical-superiority claims are made.

\begin{table}[t]
  \caption{NYC request-share tracking over 730 days. Lower values are better;
  terminal average-payoff distance and mean daily mismatch are distinct metrics.}
  \label{tab:switching-nyc}
  \centering
  \begin{tabular}{lrr}
    \toprule
    Method & $\delta_T$ & $M_T$ \\
    \midrule
    One-switch (lag base) & $9.40\cdot10^{-4}$ & 0.155910 \\
    Fast only & $6.44\cdot10^{-4}$ & 0.573865 \\
    Lag safe only & $1.67\cdot10^{-4}$ & 0.131153 \\
    Window ($W=16$) & $1.31\cdot10^{-3}$ & 0.104084 \\
    Block safe only & $3.47\cdot10^{-2}$ & 0.109093 \\
    \bottomrule
  \end{tabular}
\end{table}

\begin{figure*}[t]
  \centering
  \includegraphics[width=\textwidth]{\AAMASExperimentFigureDir/switching_nyc_dynamics.pdf}
  \caption{NYC switching dynamics. (a) Monitored residual $E_t$ and budget
  $G_T=1$; the fast-only counterfactual continues after the master's crossing.
  (b) Absolute full-target distances $\delta_t$, displaying values below
  $10^{-5}$ at that floor; no slopes are fitted.
  (c,d) Cumulative daily mismatch differences
  $\sum_{s\le t}(\|u_s^{\rm ours}\|_2-\|u_s^{b}\|_2)$ against fast-only and
  window; negative values favor the master on daily mismatch.
  The dotted line marks day 44; safe mode starts on day 45. No uncertainty bands.}
  \Description{Four panels show a real threshold crossing at day 44 on the fixed NYC demand-share trace. The master's residual stops accumulating when safe mode begins. Absolute average-payoff distances are shown for all five methods. Cumulative daily mismatch is lower for the master than fast-only and higher than window. The vertical dotted line marks the crossing; the shaded later region is safe mode.}
  \label{fig:switching-nyc}
\end{figure*}

\paragraph{Original-budget diagnostic.}
To exercise Proposition~\ref{prop:safe} unchanged, a constructed scalar game
uses $P=[0,1]$, $L=\operatorname{conv}\{0,e_1,\ldots,e_m\}$,
$u(p,\ell)=p-\sum_i\ell_i$, and $p^\star(\ell)=\sum_i\ell_i$.
For $T=16384$, 128 zero-demand rounds precede distinct unit types.
With $k=1$, $G_T=6T^{3/4}=8688.9281$; crossing occurs at $t=8817$,
leaving 7567 block-safe rounds. Terminal distances are 0.602296 for
the master, 0.992126 for fast-only, and 0.124300 for block-safe-only;
lag and window controls are below $5\cdot10^{-4}$.
The new labels have identical aggregate payoffs: this deliberately
pathological geometry tests switching execution, rather than realistic
opponent learning or minimax rates. The supplement supplies all trajectories,
separate-year NYC checks, numerical diagnostics, and the earlier frozen CAGE
study, which has no switch. Neither new test evaluates the $q=4$ exponent.
